\chapter{临床研究中的定性方法}
\chapterauthors{Daniel Dohan}
\sourcepage{454}
访谈、焦点小组及其他定性方法，能够为临床实践和研究中尚无可靠检验、调查或其他数值测量方法，或无法、难以采用这些方法的领域提供信息。它们还能整体呈现临床世界的复杂现实，保留其中丰富的社会和情境因素；这种解释是数值资料无法提供的。定性方法已被用于研究社会敏感或带有污名的疾病和处境，纳入弱势人群及研究中代表不足的群体，并将研究所得的依据用于日常实践纷繁复杂的现实（1--3）。

虽然\term{定性研究}（qualitative research）采用的是提出开放式问题、观察临床互动等常识性做法，它在卫生研究中仍存在争议（4，5）。本章讨论定性研究的基本假定、它与定量研究的区别，并举例说明定性研究的用途；还将介绍设计、实施和传播定性研究时的实际问题。

\section{什么是定性研究？}
定性研究不依赖数值测量、预先设计的检验或封闭式调查，而是帮助人们理解经历、意义和文化。研究者与参与者接触以收集资料，再解释他们的叙述性回答。以下三个例子说明研究者如何参与并塑造定性研究。
\sourcepage{455}
这些例子突出了临床定性研究的若干核心特征。每项研究都采用目的性选择而非随机选择参与者。研究团队与参与者密切合作收集资料；在肿瘤学研究中，他们甚至针对每位患者调整问题。医疗之家研究最初聚焦实施，后来随着资料收集程序根据新见解作出调整，转向研究撤除实施。EngageUC 则试图借助参与者的日常视角，了解加利福尼亚州不同社区如何看待生物样本库。

这些核心特征引出了一些棘手问题：研究者与研究如此密切地交织在一起，怎样才算无偏倚的研究？提出开放式问题的研究者，能否像封闭式调查一样做到没有偏倚？解释叙述性回答，与解释临床检验有何异同？

为回答这些问题，定性研究者首先明确考虑哪些类型的问题可以研究（11）。一些定性研究者与大多数医学研究者一样，追求客观、无偏倚的方法，采用\term{实证主义}（positivism）研究范式。实证主义者认为，宇宙受客观规律支配，科学能够认识这些规律。实证主义研究试图用精确、无偏倚的测量回答由假设驱动的问题，从而使科学家发现普遍规律。由于这种范式着重精确测量和假设检验，采用实证主义的定性研究者并不多。

一些定性研究者认为宇宙由客观规律支配，却不认为科学能够确定无疑地发现它们。\term{后实证主义}（post-positivism）者指出，科学与任何人类活动一样，总是在特定社会情境中开展，因此应审慎看待实证主义者关于测量无偏倚的主张。后实证主义者相信，研究能够帮助我们更准确地认识世界，但由假设驱动的设计并非唯一有用的工具。他们努力保持无偏倚、客观地回答研究问题，同时承认研究情境始终是理解发现的中心因素。
\sourcepage{456}
其他定性研究范式更强调情境，甚至质疑研究应以发现一套普遍规律为中心的观念。\term{建构主义}（constructivism）认为，个人和群体通过社会过程建构出客观现实的不同版本；\term{批判理论}（critical theory）研究则着重考察客观现实的性质如何反映社会中的权力关系（12）。

在每种范式下，定性研究既可运用\term{演绎}（deduction，即研究假设决定测量内容），也可运用\term{归纳}（induction，即根据观察形成假设）的逻辑。定性研究常用研究过程中出现的资料形成新想法和新假设，甚至改变研究问题、方案和场所。它不试图用数字界定变量，而是提供丰富的描述，常以叙述性报告为依据。样本通常较小，取决于研究者认为理解所研究概念、问题和过程需要多少参与者。定性研究者通过选择目的性样本、与参与者互动收集资料，来实现研究的严谨性和可重复性。定量与定性研究的重要区别见表14.1。

\begingroup\small
\begin{longtable}{L{.17\textwidth}L{.37\textwidth}L{.37\textwidth}}
\caption*{\textbf{表14.1\quad 定量研究与定性研究的比较}}\\\toprule
研究内容 & 定量研究 & 定性研究\\\midrule\endfirsthead
\toprule 研究内容 & 定量研究 & 定性研究\\\midrule\endhead
研究重点 & （见第2章）研究问题事先确定。 & 研究问题可在研究过程中改变，甚至在此时才被发现。\\
研究参与者 & （见第3章）明确界定：选择（纳入和排除）标准规定哪些人有资格参加。抽样：代表目标人群。 & 纳入和排除标准灵活，旨在包容不同参与者，可在研究过程中改变。目的性、滚雪球式抽样，可能不具有代表性。\\
\sourcepage{457}设计 & （见第8--12章）横断面研究、队列研究、病例对照研究、随机试验。 & 观察性；探索性、描述性、比较性。\\
样本量估计 & （见第5--6章）依据估计的效应量及其变异、$\alpha$ 和 $\beta$。 & 足以充分理解概念、问题和过程。\\
测量与资料 & （见第4章）程序固定，采用标准化资料收集表；结构式问卷、体格检查、实验室或影像检查、结局判定；预测变量和结局变量及协变量。 & 灵活、互动，使用现场笔记、录音或录像；民族志观察、深度访谈、焦点小组、书面或电子文件分析；记录资料收集情境和过程的备忘录。\\
数据管理 & （见第19章）标准数据库软件；编码；数据清理。 & 标准资料格式；编码方案；对资料编码。\\
分析 & 使用统计推断，检验预先计划的假设。 & 反复审阅编码结果和备忘录，识别研究发现及反例。\\\bottomrule
\end{longtable}\endgroup

\textbf{定性研究通过揭示个人经历、组织动态和健康结构性决定因素的性质，促进临床研究。}
\begin{itemize}
\item 在个人层面，定性资料对于理解患者、照护者和临床人员如何经历疾病与照护至关重要。这些理解可能是形成想法的关键：如何在多样化或弱势人群中建立或解释有关健康、疾病与照护的定量测量。
\item 在组织层面，定性方法揭示医疗服务组织复杂的内部动态，帮助临床研究者实施改善医疗服务的策略。
\sourcepage{458}
\item 在社区参与式研究中，定性方法呈现受医疗服务和临床研究影响者的视角。借此，临床研究者可以了解社区成员的健康如何反映文化和制度。这些认识可以改善研究者概念模型的效度以及假设的合理性，也能产生关于疾病原因和后果的新想法，并提出适应具体人群的健康改善干预策略与方法（见第15章）。
\end{itemize}

\subsection{定性研究的三种取向}
区分三类定性研究目标很有帮助，因为每类目标都需要略有不同的设计和方法（13）。这些术语在定性研究中的特定含义，与定量研究中的用法不同。
\begin{itemize}
\item \term{探索性研究}（exploratory studies）：所有定性项目都包含一定程度的探究，但探索性研究涉及此前没有研究过的主题，重点是描绘新的领域。研究者知道从何处开始——或许是令人困惑的临床发现，或前景良好的干预出乎意料地失败——却无法预知会发现什么。由于不知道探索将通向何方，方案最多只能是暂定的，需在研究中修订、重新聚焦和扩展。对未知主题不可能事先建立抽样框，因此研究者依赖滚雪球式资料收集，即请参与者推荐其他人（14）。同样，也没有提出分析计划的依据，因而常采用\term{扎根理论}（grounded theory）这一套归纳性资料分析程序，发掘新的理论和理解（15，16）。
\item \term{描述性研究}（descriptive studies）：这类研究更进一步。已有知识足以使研究聚焦某一领域，但尚不足以提出完整方案。当研究者想让研究工具或方案适用于新的目标人群时，常采用这种方法。他可能了解工具在一组受访者中的表现，但需探查这些经历和理解是否也适用于社会或文化特征不同的另一组人。
\sourcepage{459}
\item \term{比较性研究}（comparative studies）：其探究逻辑最接近定量研究。良好的比较性研究既阐明研究场所或参与者的相似之处，也识别重要差异。因此，它包含演绎逻辑，使研究者可以事先描述场所、资料收集程序和分析方法；但与定量研究的相似性也就止于此。定性比较研究不提出或检验假设，因为非数值资料无法进行统计学显著性检验。研究者仍接受可能改变研究问题、方案或现场的新归纳发现；这种灵活性在定量研究中很少可能实现。
\end{itemize}

\subsection{探索性、描述性与比较性取向的比较}
三类研究都可用于回答临床研究、实施科学和社区参与式研究中的问题（表14.2）。
\begingroup\small
\begin{longtable}{L{.11\textwidth}L{.26\textwidth}L{.26\textwidth}L{.26\textwidth}}
\caption*{\textbf{表14.2\quad 按设计分类的临床研究、实施科学和社区参与式研究问题示例}}\\\toprule
研究设计 & 临床研究 & 实施科学 & 社区参与式研究\\\midrule\endfirsthead
\toprule 研究设计 & 临床研究 & 实施科学 & 社区参与式研究\\\midrule\endhead
探索性 & 患者如何理解健康状况或就医经历？（见例14.1“试验知识”） & 诊所将一项干预付诸实践时，如何与之互动？ & 社区成员认为某一健康状况的哪些方面最重要，为什么？\\
\sourcepage{460}描述性 & 为提高工具的效度和信度，研究者需要了解参与者完成工具的哪些情况？ & 一个组织为何决定撤除曾热情支持的项目？（见例14.2“医疗之家”） & 在其他社区中，人们如何理解和应对某个社区的健康关切？\\
比较性 & 患者文化与组织文化如何共同塑造养老院住民的决策？ & 在一组相似诊所中，组织动态如何支持或妨碍成功实施？ & 不同社区对于如何改善卫生科学实践有哪些共识？（见例14.3 EngageUC）\\\bottomrule
\end{longtable}\endgroup

\begin{bookbox}{例14.1\quad 试验知识：探索癌症患者对临床试验的认识}
\term{民族志}（ethnography）是记录群体、团队、组织或社区社会互动和行为的定性社会科学方法，通过收集定性资料，对一种文化作细致、整体的分析（6）。例如，一项民族志研究对78名晚期癌症患者进行深度访谈，探究肿瘤科医生如何招募患者参加早期临床研究（7）。研究团队首先观察肿瘤门诊的诊疗，与负责招募患者的临床医生交谈。他们了解到，试验招募的非正式环节，早在医生请患者签署知情同意书之前就已开始。肿瘤科医生说，他们甚至在讨论参加研究之前，就会考虑谁是“好的研究参与者”（8）。由此产生了几个问题，例如患者是否知道这一过程，以及被邀请前对临床试验知道多少。
\sourcepage{461}
团队据此编制访谈方案。为避免影响回答，除非患者自己提起，访谈者不会讨论参加试验的话题。大多数患者在首次访谈中便主动谈到这一点。许多人描述了疾病较早阶段参加研究的经历。有些人对此十分热情，并盼望以后再参加；另一些人的感受较为复杂，或担心自己没有资格参加。团队还发现，患者自认为对试验了解的程度，与其知识的准确程度往往不相称。熟悉甚至精通肿瘤学技术用语，并不等于实质上理解临床试验。

这些发现提醒我们，患者说出的内容与真正理解的内容并不相同。患者可能努力表现得知识丰富，也许是想呈现自己是“好的研究参与者”，从而获得参加试验的机会。但照单全收他们的话，可能反而不利于患者和临床试验过程。
\end{bookbox}

\begin{bookbox}{例14.2\quad 撤除医疗之家：描述一种新医疗服务模式为何被取消}
\term{关键知情人访谈}（key informant interviews）是了解专家观点与视角的定性工具。访谈通常为开放式、半结构式：访谈者围绕预定主题自行组织问题并决定何时提问，受访者用自己的语言回答。与所有定性研究一样，当出现有趣却未预料的问题时，访谈者可以自由提出新问题继续探查。这种方法尤其适合了解专家对其熟知主题的看法，以及探索领导者如何作出影响深远的决定。
\sourcepage{462}
对38名领导者、管理者、临床人员和工作人员的关键知情人访谈，帮助解释了某大型医疗集团为何放弃“以患者为中心的医疗之家”（Patient Centered Medical Home）这一新政策举措（9）。对政策制定者而言，该决定提出了令人不安的问题：为什么医疗之家没有实现鼓励、激励和支持以患者为中心的照护的目标？

研究发现，组织内的利益相关者从未把以患者为中心的照护与医疗之家政策联系起来。即使任由医疗之家认证到期，临床人员、领导者和工作人员仍认为机构致力于落实以患者为中心的照护。组织内部，一组管理者负责保证患者中心照护，另一组负责争取正式的医疗之家认可。资深医生和领导者认为，患者中心照护始终是组织优先事项，但负责认证的管理者没有把其工作与这一核心价值联系起来。他们认为，组织寻求认证是因为一个大型雇主团体提供了奖励。到续认证时，奖励已终止，管理者便优先处理其他工作，让认证失效。政策制定者确实有理由担心：医疗之家政策从未在组织内取得认同；但没有续认证，并不表示组织不再优先提供以患者为中心的初级卫生保健。
\end{bookbox}

\begin{bookbox}{例14.3\quad EngageUC：比较不同社区对生物样本库研究的看法}
定性研究者组织4--10人的\term{焦点小组}（focus group），了解社区对某一主题的态度。小组回答主持人提出的一组问题。（通常有两位主持人，一人推动讨论，另一人留意非语言互动，并确保每个人都能发言。）加利福尼亚大学（University of California，UC）想扩充研究用生物样本库，使其更充分地纳入州内多样化人群，于是在州的南北两地组织社区成员焦点小组，了解看法并征求建议（10）。
\sourcepage{463}
项目旨在寻找让 UC 更好回应州内各人群需求的方法。团队从具有不同种族、族裔历史和传统的数十个社区中，有目的地选择51名社区成员，呈现加州的多样性。焦点小组分别用英语和西班牙语开展。领域专家讲解样本库运行的技术内容，经验丰富的主持人引导社区成员讨论。

同时访谈一组人看似比个别访谈高效，但焦点小组也有复杂之处：团队需要广泛出行，并安排参与者的交通、膳食和补偿；还要处理大量定性资料录音、转录、翻译和分析的技术问题。最后一天，参与者向 UC 样本库领导者提出建议并表决，形成23项共识建议，另有24项获得支持但未达成共识。

焦点小组帮助研究团队和样本库领导者了解加州人如何看待大学及州内的样本库研究。来自不同背景的普通公众在一些重要领域达成共识：应广泛开展样本库教育，自由分享信息和数据应得到大力支持，研究者应接受更多监督并承担责任，知情同意应超越法律最低要求，与提供样本的患者充分交流。
\end{bookbox}

\sourcepage{464}
探索性研究有一些共同特点，常采用建构主义范式。问题宽泛而开放，反映描述和发现的目标，不暗示研究者预期在概念之间发现某种特定的因果或非因果联系。

描述性定性研究常在新的目标人群中考察已有发现，或考察为某一人群开发的工具在另一人群中的接受情况。它可以探索一个社区如何理解在别处出现的健康关切。在实施科学中，描述性研究可以考察临床结局不同的诊所（例如糖尿病控制情况不同者）中的医患沟通。

比较性设计预先考虑哪些群体值得比较，以及可能存在什么差异。设计常由理论驱动并采用目的性抽样，例如在随机试验的干预组和对照组中选择部分成员进行定性比较。其发现有助于解释不同群体的行为、信念和结局模式。

\subsection{定性方法与定量技术的比较}
所有定性研究都着重解释健康与医疗服务，例如患者如何选择治疗、诊所如何应对循证实践的变革，或某社区为何承受某种疾病不成比例的负担。研究者在日常生活中观察这些问题，补充定量方法的发现。定性研究常提供个人层面的资料，帮助解释群体层面的发现。无论作为同时包含定性和定量方法的\term{混合方法研究}（mixed-methods study），还是大型项目中独立的探索性部分，定性资料都常是定量研究的有益补充。相对少量的定性资料即可提供见解，提高定量研究的效度，或使主体研究避开未预料的困难。

\sourcepage{465}
\section{何时采用定性方法}
定性资料有助于开发测量方法、把握复杂概念和探索新想法。

\subsection{开发新测量方法或改进已有方法}
定性方法有助于开发封闭式问题。在研究早期，半结构式访谈可识别潜在参与者认为某个主题的哪些方面重要。访谈和焦点小组能帮助判断条目对不同受访者是否适当，揭示可能的误解，并指导完善最终条目。它们也能帮助解释定量结果。例如，分析日常文化的民族志资料能够解释干预为何产生某种结果。由于开发和完善测量以已有知识为基础，旨在深化对某一现象理解的描述性设计往往适用。

\subsection{把握难以定量测量的复杂概念}
临床研究者关心的某些概念并不适合定量测量，主要涉及三个领域：
\begin{itemize}
\item 在个人层面，研究者可能关心行为对参与者的意义。例如，患者决策研究可借助定性资料超越决定本身，探索患者如何作出决定。
\item 很难用定量测量把握双人互动的细微差别，例如医患之间或患者与家庭照护者之间的互动。定性资料常能揭示这种关系的动态，尤其是从关系双方的视角收集并相互参照印证的资料。
\sourcepage{466}
\item 诊所、社区、医疗服务组织等复杂系统会产生并维系自身文化。民族志和其他定性测量能够把握文化的丰富与微妙之处，因此被视为这类研究的金标准。
\end{itemize}

\subsection{探索新想法}
定性方法能够发现临床医学的新侧面，开创性研究包括《白衣男生》（\textit{Boys in White}，考察医学院如何通过社会化使学生形成医生角色）（17）、《临终意识》（\textit{Awareness of Dying}，考察患者、照护者和临床人员如何面对医院中的死亡经历）（18），以及《宽恕与铭记》（\textit{Forgive and Remember}，考察外科医生如何处理错误）（19）。事实上，临床人员在照护患者、教学和领导卫生系统时，可以运用观察、记笔记、提出开放式问题并追问回答含义等定性工具。这些活动可使他们进一步认识患者、学生和同事的经历，进而提高开展研究的周密程度、可重复性和影响。

\section{使设计适合研究问题}
定量研究的许多惯例，如随机化和盲法的价值、测量并调整混杂因素的需要、估计样本量的重要性，在定性研究中并无现成对应。在缺少这些指引时，研究者必须反思并预想研究问题、已有的场所和参与者知识、预期读者以及研究可能如何展开，据此设计研究。

探索性研究考察新颖和未知的事物，因而需着重考虑新颖之处为何重要以及如何重要。设计伊始，研究者就应考虑成果的读者。设想一项研究探索弱势人群未获满足的医疗需求。一类读者可能已熟悉问题并了解其重要性；针对他们，研究应收集识别新问题的资料，以归纳逻辑管理和分析资料。
\sourcepage{467}
可以先深度访谈少数关键知情人，再按最初访谈的发现扩大范围。研究者可着重通过团队会议共同审阅资料，并撰写\term{分析备忘录}（analytical memos）：用短文指出新出现的兴趣主题，指导持续收集资料。

同一个问题，如果旨在为日后改变临床实践的研究提供初步资料，则可能需要不同的探索性设计。它必须说服可能对小型定性研究持怀疑态度的临床读者，因此很可能需要更多参与者，即使每人收集的资料更少。开放式调查问题或结构化的开放式访谈方案，加上周密的数据管理和分析，可能比对一家诊所进行深度民族志研究更合适。

探索此前未研究主题时，设计应着重考虑如何进入现场、接触参与者。研究者可能需要借助已有关系，并说明如何随着研究推进与参与者建立联系。探索性设计关注研究过程，也要说明如何决定何时停止收集资料。研究常受组织安排或经费限制；此时应承认限制，并说明现有时间为何足以实现目标。其他情况下，研究持续到\term{主题饱和}或\term{资料饱和}（thematic or data saturation），大致指继续收集资料已难以产生新认识的时点。

比较性设计既需收集足以支持跨场所或跨参与者分析的资料，又需考虑如何接触足够多的人。为兼顾两者，研究者应采用目的性抽样策略，先建立概念模型，说明各场所或群体如何满足分析需要。据此考虑每个场所的结果可能如何影响项目发现，并在不同情况下调整资料收集和分析。
\sourcepage{468}
例如，比较研究者可能想了解高质量与低质量诊所的照护过程（如糖尿病和高血压控制）有何不同，最终借鉴高质量诊所的经验。研究者可能根据质量的定量指标，从两类中各选一家诊所开始收集资料。但如果这些指标不能很好反映诊所运行的整体状况呢？所谓高质量诊所，也许只是服务的患者条件更优越，或记录质量资料时更有技巧，并非提供照护更好。这表明起初的概念模型有误。此前的探索性或描述性研究本应提示：社会优势会混杂质量测量，诊所也可能“钻”测量系统的空子。这样便能选择照护过程确有差别的诊所，例如处于不同医疗系统或具有不同领导结构者。将这些见解纳入模型，才能更好实现识别并移用优质照护过程的目标。

\section{资料收集}
选定设计后，研究者要决定收集资料的组织安排和具体操作。

\subsection{确定样本量}
“样本量多大合适？”是关于定性研究最常见的问题。其目的并非获得足以检出某个效应量的检验效能，而是理解所选概念、问题和过程。许多定性社会科学研究者认为，由一名研究者开展的适度规模研究，可从30--40人起步。
\sourcepage{469}
最终，合适的样本量取决于问题、预期读者和组织安排，良好的样本需兼顾三者。
\begin{itemize}
\item \textbf{问题：}探索性问题通常采用小样本；扎根理论分析的经验是需要20--25次访谈。描述性和比较性研究从多个场所或人群招募，故样本通常较大。资料收集与分析反复交替，使招募可随新问题调整。研究者逐步深入理解问题时，可能第一处访谈15人、第二处10人、第三处5人便已足够；回访进行确认性访谈也许再需10人，总数仍回到30--40人这个“恰到好处”的估计。
\item \textbf{读者：}定性研究的重要目标是说服政策制定者、临床人员、研究者和期刊编辑相信发现有效。这些读者习惯大样本，可能担心小型选定样本。增大样本或使各组人数相近，未必是实现定性目标所必需，却可能回应这些顾虑，增加发表机会。
\item \textbf{组织安排：}定性研究需要大量人力。获得并保持进入场所、接触参与者的机会需要时间。许多研究由一人承担；即使是团队项目，主要研究者仍负责与现场建立和维系关系。研究者的时间常是限制因素，不宜把样本量增加到超出科学需要的程度。
\end{itemize}

\subsection{访谈还是焦点小组？}
另一个常见问题是采用访谈还是焦点小组。焦点小组一次可访谈6--8人，似乎更高效。
\sourcepage{470}
但两者收集的是不同类型的定性资料，应根据研究问题选择，而非根据效率和规模的印象。访谈揭示个人观点、信念和经历；焦点小组揭示成员彼此之间以及对主持人的反应。

若想了解不同社区或群体中个人的观点，一系列个别访谈通常更合适。焦点小组有其内部动态，例如话多者可能塑造集体观点。熟练主持人虽能部分克服这种倾向，焦点小组适合揭示群体互动，却不是探索个人观点的工具。例14.3的样本库研究旨在征求社区建议，因此采用焦点小组，观察不同社区的个人如何相互讨论并形成建议，正是合适的方法。

\subsection{管理研究现场与参与者}
定性访谈需要参与者和研究者投入大量时间。收集资料的工作常由主要研究者承担，他可能就是团队唯一成员。深度访谈研究也不像采集生物样本或发放封闭式问卷那样容易分工。每次访谈的展开都略有不同，访谈者须能熟练、灵活地运用半结构式方案。即使研究规模足以组建多人团队，负责人仍需大量亲自参与现场人员的培训和监督。

\subsection{保护保密性}
定性资料给匿名化与保密带来独特困难。有些参与者仅凭其描述的职位或故事就很容易被辨认。
\sourcepage{471}
研究者可能询问敏感问题，例如患者是否遭遇不平等待遇，或质量改进项目为何失败。团队也可能需要制定程序，使其能在不泄密的前提下，用其他访谈交叉核查一次访谈所得的认识。

研究者必须保护参与者的保密性。大多数人知道不应利用权力有限的参与者，但有效的定性研究也可能给拥有权力和优越地位者带来隐蔽风险。研究临床医生或医疗组织领导者时，熟练研究者会与这些知情人建立信任，使他们谈及平时保密、可能受污名的行为或态度。如果资料公开，其声誉可能受到严重影响。

\subsection{访谈提纲}
制定\term{访谈提纲}（interview guide）是改进执行、预见和避开困难的机会。研究者通常先拟定初稿，再在整个项目中不断修订。提纲既确保覆盖主要问题，也为意外主题留出空间。

许多研究适合采用\term{主问题与追问}（stem and probe）方式设计访谈和焦点小组问题。主问题引出核心主题，再用问题、提示或陈述追问，促成充分、周到的回答。例如，在癌症患者试验知识研究（例14.1）中，一个有用的主问题是：“请告诉我，你上次看肿瘤科医生时发生了什么。”接着追问患者听到医生如何解释诊断、预后和治疗选择，是否有机会提出自己的关切，以及医生是否听取并回应。
\begin{itemize}
\item 探索性提纲通常只有少数需进一步了解的主题（10个或更少）。主问题不必是正式问句，因为访谈理想上像对话。提纲提醒访谈者覆盖有关主题，但访谈往往依照自身逻辑展开，而非预设结构。
\sourcepage{472}
\item 描述性访谈更有结构。提纲可列出研究者优先关注问题的主问，再用追问促使受访者展开。
\item 比较性研究通常采用结构化提纲，着重询问受访者的经历与其他场所相似或不同之处。提纲可规定各问题如何提问，以帮助聚焦分析。
\end{itemize}

进入现场前，提纲向潜在资助者和合作者展示程序。进入现场时，研究者用它试访并培训访谈者。在这些过程中，提纲不断改进：调整长度，修订不清楚的问题，改变顺序。正式开展后，还可根据新发现进一步修订，删去重复或无意义的问题。最后版本用于论文写作和修订时描述研究。

\subsection{收集资料}
定量策略力求样本代表目标人群，并用固定程序消除资料收集偏倚（表14.3）。相比之下，定性研究者选择目的性样本并积极参与资料收集，以实现严谨性与可重复性。代表性和偏倚是临床研究熟悉的原则，目的性和参与性则较陌生。目的性抽样要求研究者考虑所选场所和参与者能提供何种信息以回答问题；灵活的收集策略则使其发现挑战既有理论的意外结果。
\sourcepage{473}
\begin{bookbox}{表14.3\quad 定量与定性研究中严谨且可重复的资料收集}
\begin{tabularx}{\linewidth}{L{.17\linewidth}YY}\toprule
研究活动 & 定量研究 & 定性研究\\\midrule
选择参与者 & 代表性：从具有代表性的群体中随机抽样。 & 目的性：根据理论需要选择场所和参与者。\\
获得资料 & 限制偏倚：固定收集程序，不受研究者影响。 & 促进参与：灵活收集，包含研究者与参与者的互动。\\\bottomrule
\end{tabularx}
\end{bookbox}
例如，EngageUC（例14.3）希望纳入以往未充分参与基因组研究讨论的西班牙语受访者。因此需全程安排同声传译；团队也讨论了分组讨论应采用双语加传译，还是单一语言而无需传译。最终，参与者表示希望采用能自由交流的单语分组，帮助作出了决定。

\section{资料管理、分析与报告}
新手与有经验的研究者都可能觉得定性资料管理和分析有些神秘。习惯用统计软件清理、分析数字的人，面对等待分析的访谈转录稿可能无所适从。资深研究者也说，沉浸于定性资料并形成解释性分析的思考过程很难描述。

\term{定性资料分析}（qualitative data analysis，QDA）软件帮助管理文本，却不真正代替人分析；它帮助分析者组织资料（即使适度规模的研究也可能有数百页资料），并跟踪不断出现的结果。
\sourcepage{474}
\subsection{定性分析的目标}
定量与定性研究均通过分析完成两项关键任务：避免错误发现，并以读者认为可信的方式交流发现（表14.4）。\textbf{定性研究者反复审阅、探究资料，以完善结果并减少解释错误。}这一过程从收集资料开始，贯穿整个项目。关键步骤是建立\term{编码}（codes），把握分析中需考察的分析性与概念性问题。每个项目都有独特的一组编码，其数目和定义反映项目规模、目标、研究范式，以及演绎与归纳逻辑所占的程度。一般来说，规模较小（实证资料较少或概念目标较有限）的项目，以及后实证主义传统下的演绎项目，编码较少。完整编码包括：明确的定义；纳入与排除标准，说明它与相关编码的相似和区别；适用资料的例子（必要或有帮助时也列出不适用的例子）；以及\term{资料备忘录}（data memo），描述编码的来源及其在反复收集、分析中的发展和使用历史。

\begin{bookbox}{表14.4\quad 定量与定性研究中严谨且可重复的资料分析}
\begin{tabularx}{\linewidth}{L{.17\linewidth}YY}\toprule
研究活动 & 定量研究 & 定性研究\\\midrule
避免偶然发现 & 统计：用统计推断检验假设。 & 过程：反复审阅资料（常包括编码和备忘录），完善结果。\\
\sourcepage{475}使他人信服 & 透明性：以标准、熟悉的格式呈现资料。 & 情境：提供足够的叙述情境，使读者能够解释结果。\\\bottomrule
\end{tabularx}
\end{bookbox}

研究者把编码组织为\term{编码手册}（codebook）。与定量数据集的变量代码簿类似，它保存编码定义、示例及其他信息，也展示编码之间的关系，通常按层级把具体编码归入较大分析类别，常画成树状结构。编码手册使概念模型得到操作化；模型的中心要素应反映在其结构中。随着新编码被发现和定义、团队认识新的概念关系，手册也会演变。达到主题或资料饱和时，手册定稿：此后不应再需增加编码，概念关系也应足以支持发布最终手册，对全部资料作确定的编码。这是定性研究的重要转折点。

有了编码和手册，团队即可集中编码资料。有些编码识别主要主题，有些用于分类和整理，例如标明所观察互动发生的地点。团队探索编码资料中的模式时，这类组织性编码尤其有用。团队以分析备忘录记录模式。（较早阶段，资料备忘录记录的是编码演变。）此时，分析备忘录记录研究发现的形成过程，把握分析者新出现的认识和想法。其中有些想法最终没有结果：新增资料可能表明最初观察只是一次性事件，或与其他资料矛盾；团队讨论可能认为它应并入其他发现，或只是某一分析者个人的解释，尚不能获得团队认同。
\sourcepage{476}
另一些备忘录则把握最终形成研究发现的核心想法。例如，“试验知识”项目（例14.1）的一份备忘录指出一个反差：一名患者熟知试验药物许多生僻技术细节，却似乎不理解临床研究，尤其早期临床研究的基本目标。其他分析者认同这一解释，随后检查编码资料，考察其他患者是否也表现出类似知识模式。这一认识最终成为关于晚期癌症患者理解试验过程的主要发现之一。

定性研究产生的文本和叙述资料，不易纳入旨在提供透明性、促进可重复性的标准报告格式。因此，分析者需提供情境细节，展示自己的解释和发现，使他人信服结果的效度。

\subsection{使用定性资料分析软件}
计算机辅助 QDA 已很普遍，但没有改变管理定性资料的基本过程。以前，分析者通过阅读、标记、分类来理解资料，并写备忘录记录思考、完善发现。这些步骤仍是避免武断发现、说服读者的必要工作；可以在纸质访谈稿页边用彩笔批注，把转录稿复印后分堆，甚至用 Word、Excel 等通用办公软件完成。

QDA 软件主要是带有批注、排序和筛选功能的定性资料数据库。多数软件以扎根理论传统为基础，旨在帮助探索性编码及主题内容分析。
\sourcepage{477}
在具有这些目标的探索性项目中，它可有效用于：
\begin{itemize}
\item 多名成员参与编码时，通过版本控制保持资料完整性；
\item 便于计算编码者间信度的定量得分，记录编码的效度；
\item 创建和维护探索性研究产生的主题编码，特别是管理和分析中出现大量编码时；
\item 生成单个编码或编码组的报告，包括导出特定编码、或编码手册特定部分的原始资料。
\end{itemize}

\subsection{管理和分析定性资料的过程}
\term{资料管理方案}（data management protocol）说明访谈、焦点小组及其他现场资料如何转为一致的分析格式。首先描述研究者如何进入现场、接触受访者、获得并记录知情同意，以及如何用数字录音设备（最好两台，以免一台故障导致丢失）或笔记收集叙述资料。标准化处理包括转录访谈、说明如何识别焦点小组中的发言者，以及如何把笔记本或录音中的现场记录转成文字处理文件。

建立 QDA 数据库还需决定文件如何组织和命名。初始手册应包含描述性或比较性研究的演绎问题，也可列出探索性研究预期发现的主题。但所有项目的手册都会随时间演变。确定最终编码方案常是瓶颈。现场工作期间，团队可多次开会审阅资料、讨论编码，以完善分析概念并纳入意外发现。
\sourcepage{478}
阅读、讨论、再阅读不可或缺；纸笔编码和用白板追踪概念、编码及想法也会有帮助。

有了资料管理方案和编码手册，研究者可使用软件或手工开展分析的中心工作。关键决定是侧重归纳还是演绎。归纳分析在探索性研究中较常见，如扎根理论：对资料作归纳编码以发现主题，再把主题组织起来，提出新理论。

比较性或描述性研究至少部分采用演绎分析，先列出拟考察的一组概念或主题，类似定量分析的变量。例如，可采用\term{框架分析}（framework analysis）：每行为一个研究概念，每个个案构成一列，以比较个案。单元格填入能说明该概念的访谈引文或其他定性资料（20）。

研究者应判断研究更适合归纳还是演绎，选择相应策略。应避免用扎根理论等归纳方法分析具有演绎性质的研究资料，例如分析者已经能明确说出预期发现的研究。此时，归纳分析可能错误地暗示新认识“证实”了研究者的预期。更站得住脚的做法是承认访谈提纲或其他设计与程序可能诱导受访者作出确认性回答；采用更注重发现、考察矛盾资料的演绎方法，更可能产生有说服力的发现。咨询导师和受过定性训练的研究者，有助于识别并避免这类错误。

\subsection{撰写和分享结果}
\sourcepage{479}
资料管理方案和记录团队思考的备忘录，可成为论文方法部分的初稿；资料管理与分析过程构成结果部分的中心。
\begin{itemize}
\item 探索性研究的结果包括归纳发现的主题。发现往往多于能够报告的内容，应优先选择与预期读者最相关者。选择最有说服力的示例并无不妥，但若存在矛盾结果，也必须提供，以保证解释准确。
\item 描述性设计应报告符合预期主题的发现和意外结果。归纳发现还可讨论研究如何根据新资料的分析而演变。
\item 比较性研究主要报告预先探问的场所之间的异同，尤其能说明主要主题者，也应包括意外发现。
\end{itemize}

定性研究者常向利益相关者分享发现，包括参与者及以其他方式帮助研究的人。参与者往往想知道研究者如何理解其故事和见解，以及其经历与别人如何异同。分享时必须保护身份：对研究者而言平淡无害的发现，在利益相关者之间可能引发争议；个别访谈或焦点小组中分享的观点，公开后也可能引起担忧。

在可预见的将来，定性研究仍只占已发表临床研究的一小部分。大多数期刊限制篇幅，使作者难以充分展现令结果可信的叙述丰富性和情境特殊性；编辑和审稿人也可能缺乏相关接触与专业知识。
\sourcepage{480}
写稿时灵活处理可能有助于克服这些障碍。例如，有些期刊不把表格和框内文字计入字数，可考虑将引文和支持资料放在其中。应仔细选择表头：探索性研究适合用引文展示相关主题群，描述性和比较性研究可按场所或是否符合预期组织资料。选择投稿期刊时，可以先向编辑询问；编辑也可能欢迎推荐合格审稿人。

最后，向高影响力期刊投稿需注意。临床研究者常先投最高等级期刊，不成功再往下投。但定性论文改投面向不同读者的期刊时，可能需要大量重写，尤其是字数上限不同者。压缩本已精练的文章，可能费时又痛苦。首次就投向更愿意接受定性研究的期刊，通常更好。

\section{小结}
\begin{enumerate}
\item 定性方法使研究者能整体考察复杂、依赖情境或敏感的主题。
\item 定性研究用于探索新现象，描述已有概念或工具在新场所、新人群中的应用，以及比较复杂社会过程和组织。
\item 定性研究要求自我反思，首先要选择适当的探索性、描述性或比较性设计。
\item 收集资料时，应根据问题选择场所和参与者，以灵活策略促进参与，确保获得丰富资料。
\sourcepage{481}
\item 定性研究反复交替收集、审阅与分析资料，并在报告中提供情境细节。
\end{enumerate}

\begin{references}
\item Hoff TJ, Witt LC. Exploring the use of qualitative methods in published health services and management research. Med Care Res Rev. 2000;57(2):139-160.
\item Mays N, Pope C. Rigour and qualitative research. BMJ. 1995;311(6997):109-112.
\item Weiner BJ, Amick HR, Lund JL, Lee SYD, Hoff TJ. Review: use of qualitative methods in published health services and management research: a 10-year review. Med Care Res Rev. 2011;68(1):3-33.
\item Devers KJ. How will we know “good” qualitative research when we see it? Beginning the dialogue in health services research. Health Serv Res. 1999;34(5):1153-1188.
\item Greenhalgh T, Annandale E, Ashcroft R, et al. An open letter to the BMJ editors on qualitative research. BMJ. 2016;352:i563.
\item Reeves S, Kuper A, Hodges BD. Qualitative research methodologies: ethnography. BMJ. 2008;337:a1020.
\item Garrett SB, Koenig CJ, Trupin L, et al. What advanced cancer patients with limited treatment options know about clinical research: a qualitative study. Support Care Cancer. 2017;25(10):3235-3242.
\item Joseph G, Dohan D. Diversity of participants in clinical trials in an academic medical center: the role of the “good study patient?” Cancer. 2009;115(3):608-615.
\item Dohan D, McCuistion MH, Frosch DL, Hung DY, Tai-Seale M. Recognition as a patient-centered medical home: fundamental or incidental? Ann Fam Med. 2013;11(Suppl\_1):S14-S18.
\item Dry SM., Garrett SB, Koenig BA, et al. Community recommendations on Biobank Governance: results from a deliberative community engagement in California. PLOS One. 2017;12(2):e0172582.
\item Teherani A, Martimianakis T, Stenfors-Hayes T, Wadhwa A, Varpio L. Choosing a qualitative research approach. J Grad Med Educ. 2015;7(4):669-670.
\item Brown MEL, Dueñas AN. A medical science educator’s guide to selecting a research paradigm: building a basis for better research. Med Sci Educ. 2020;30(1):545-553.
\sourcepage{482}
\item Rendle KA, Abramson CM, Garrett SB, Halley MC, Dohan D. Beyond exploratory: a tailored framework for designing and assessing qualitative health research. BMJ Open. 2019;9(8):e030123.
\item Goodman LA. Snowball sampling. Ann Math Stat. 1961;32:148-170.
\item Charmaz K. Constructing Grounded Theory. Sage Publications; 2014.
\item Glaser BG, Strauss AL. The Discovery of Grounded Theory: Strategies for Qualitative Research. Aldine Publishing Company; 1967.
\item Becker HS, Geer B, Hughes EC, Strauss AL. Boys in White; Student Culture in Medical School. University of Chicago Press; 1961.
\item Glaser BG, Strauss AL. Awareness of Dying. Aldine Pub. Co.; 1965.
\item Bosk C. Forgive and Remember: Managing Medical Failure. University of Chicago Press; 1979.
\item Gale NK, Heath G, Cameron E, Rashid S, Redwood S. Using the framework method for the analysis of qualitative data in multi-disciplinary health research. BMC Med Res Methodol. 2013;13(1):117.
\end{references}

\chapterappendix{14A}{第14章习题：临床研究中的定性方法}
\sourcepage{483}
\begin{enumerate}
\item 在历史与当代歧视的背景下，已有充分记录表明，美国黑人对医学及医学研究的信任有限。一名研究者计划用定性研究进一步了解，不同黑人群体的不信任性质与经历可能有何不同。研究将纳入近期来自加纳的移民、来自多米尼加共和国的第二代和第三代非洲裔加勒比人，以及在美国生活了多代的非裔美国人。请提出适当的定性设计（探索性、描述性或比较性），并说明理由。
\item 例14.3的 EngageUC 用焦点小组征求建议，探讨 UC 如何使不同加州社区更好参与样本库研究。请选择使用焦点小组的最佳理由：
\begin{enumerate}
\item 考虑到加州的规模，焦点小组使研究者能高效纳入大量参与者。
\item 焦点小组使研究者能考察并理解，不同社区成员如何彼此交流、理解样本库。
\item 在群体情境中，参与者感到更安全，因此可以更自由地发言。
\end{enumerate}
\item 大多数定性研究同时包含归纳和演绎分析。考虑例14.1“试验知识”和例14.2“撤除医疗之家”：前者最适合哪种分析，后者呢？哪项研究同时包含两种分析？请解释。
\end{enumerate}
